\chapter{Scope, Method, and Claim Discipline}
\label{chap:scope}

\section{The problem being addressed}

The Riemann Hypothesis states that every non-trivial zero $\rho$ of the Riemann zeta function has real part $1/2$ \citep{riemann1859,clay2026}. It is important to distinguish this global assertion from the much simpler observation that the completed zeta function possesses a symmetry centred at $1/2$. A symmetry centre does not generally force every zero to be fixed by the symmetry. It may only force zeros to occur in symmetric families.

The present project asks whether the axis $\Re s=1/2$ can be characterized simultaneously as:
\begin{enumerate}
  \item the fixed set of the anti-linear zeta involution;
  \item the unique point of binary informational balance;
  \item the unique equal-gain two-channel state permitting exact destructive cancellation;
  \item the unique value in $(0,1)$ producing the spinorial sign $-1$ under the phase law $e^{2\pi i\sigma}$.
\end{enumerate}
Each item can be proved exactly once the relevant definitions are fixed. The difficult question is whether the zeros of $\xi(s)$ admit a canonical representation in which these structures are not merely parallel but are mathematically the same event.

\section{Variables and domains}

Throughout the monograph,
\[
 s=\sigma+it,\qquad \sigma,t\in\R,
\]
and the open critical strip is
\[
 \Sstrip=\{s\in\C:0<\Re s<1\}.
\]
The critical line is
\[
 \Crit=\{s\in\C:\Re s=1/2\}.
\]
The word ``binary'' refers to a normalized pair of non-negative weights $(\sigma,1-\sigma)$. No claim is made that the zeta function is literally a two-state physical system. Rather, the ansatz asks whether a two-channel representation can be constructed canonically from the analytic data.

The natural logarithm is denoted by $\ln$. Entropies are therefore measured in nats, so that the balanced binary entropy is $\ln2$. The complex power $2^s$ means $\exp(s\ln2)$, with the real logarithm of $2$.

\section{Autonomy and non-circularity}

A construction intended to illuminate the Riemann Hypothesis must not encode the answer in its definitions. There are several ways circularity can enter:
\begin{itemize}
  \item defining a state only at already known zeros and choosing its weights to be $1/2$;
  \item using a phase function that is set to $\pi$ precisely when $\xi(s)=0$;
  \item dividing by $\xi(s)$ and then interpreting a singular normalization as a zero detector;
  \item fitting parameters to a finite list of tabulated zeros;
  \item invoking self-adjointness without constructing a domain, operator, and spectral equivalence.
\end{itemize}
The canonical bridge demanded later in this text is designed to exclude these manoeuvres.

\begin{definition}[Canonical for the present programme]
A map, operator, or kernel is called \emph{canonical} if it is determined by previously specified zeta, theta, Hilbert-space, or symmetry data; is defined on an open domain rather than only on the zero set; is covariant under the relevant involutions; contains no parameters fitted to zero locations; and admits an independently checkable rule for its normalization.
\end{definition}

This definition is intentionally demanding. A non-canonical map can always be manufactured to send a prescribed set to any desired geometric locus. Such manufacture has no explanatory power.

\section{Claim classes}

The core repository begins with five claims. Three are elementary lemmas:
\begin{align*}
\text{SOH-L001:}&\quad \Hbin(\sigma)\text{ is uniquely maximal at }\sigma=1/2,\\
\text{SOH-L002:}&\quad \Acomp(\sigma,\phi)=0\iff \sigma=1/2,\ \phi\equiv\pi\pmod{2\pi},\\
\text{SOH-L003:}&\quad \operatorname{Fix}(\J)=\Crit,
\end{align*}
where $\J(s)=1-\overline{s}$. One claim is the open bridge:
\begin{quote}
\textbf{SOH-C001.} A canonical zeta-state map converts every non-trivial zero into exact complementary cancellation.
\end{quote}
The final claim is conditional:
\begin{quote}
\textbf{SOH-T001.} Under SOH-C001 and the specified normalization, all non-trivial zeros lie on $\Crit$.
\end{quote}

The monograph strengthens the exact layer with several additional results. In particular, the entropy deficit is a Kullback--Leibler divergence, equal-gain cancellation has a coercive quadratic expansion, the binary-spinor model admits an equivalence theorem, and unequal channel metrics move the balance point away from $1/2$. These additions sharpen both the positive structure and the hidden assumptions.

\section{What would count as success?}

There are several levels of success, and they should not be confused.

\paragraph{Level 1: structural explanation.}
One proves that complementarity, entropy, equal-gain cancellation, spinorial sign, and zeta symmetry all single out the same half-coordinate. This is achieved in the exact model developed here.

\paragraph{Level 2: canonical zero representation.}
One constructs a state or kernel from classical zeta data such that its vanishing is equivalent to $\xi(s)=0$ and its channel norms are linked to $(\sigma,1-\sigma)$. This is not achieved here.

\paragraph{Level 3: positivity or self-adjointness.}
One proves that the canonical representation is governed by a positive kernel or a self-adjoint operator, forcing the centred spectral variable to be real. This would imply the Riemann Hypothesis.

\paragraph{Level 4: accepted proof.}
Every analytic detail, domain issue, convergence statement, and equivalence is verified independently and accepted through ordinary mathematical publication and review. Nothing in this monograph claims that level.

\section{Why pursue a conditional theorem?}

A conditional theorem is useful when it compresses the unresolved problem into a precise construction target. The Hilbert--P\'olya strategy is valuable for exactly this reason: if the zero ordinates are the spectrum of a self-adjoint operator, they are real, and the Riemann Hypothesis follows. The hard part is not the final spectral implication but the operator construction \citep{berrykeating1999,connes1999,dlmf2026}.

The information--spinor theorem developed here plays an analogous role. It does not replace Hilbert--P\'olya; it refines what a two-channel geometric realization would need to accomplish. Its value is measured by whether it rules out vague analogies and generates testable mathematical subproblems.

\status{This chapter is methodological. Its definitions govern how later statements are classified; it contains no claim that RH is proved.}
